% AM_default_template.tex
%
% Editable LaTeX template for ParrotFeedback PDF reports. Style this
% however you like - fonts, colors, spacing, add a real logo image,
% whatever you want. ParrotFeedback only touches the placeholder
% keywords used below (see README.md for the full list and what each
% one is replaced with); everything else in this file is yours to edit.
%
% IMPORTANT if you add more placeholders or copy this file: don't
% mention a placeholder's exact spelling anywhere except where you want
% its value inserted - the replacement is a plain find-and-replace over
% the whole file, comments included.

\documentclass[11pt,a4paper]{article}
\usepackage[a4paper,left=2.2cm,right=2.2cm,top=2.5cm,bottom=2.5cm]{geometry}
\usepackage[utf8]{inputenc}
\usepackage[T1]{fontenc}
\usepackage{helvet}
\renewcommand{\familydefault}{\sfdefault}
\usepackage{xcolor}
\usepackage{longtable}
\usepackage{array}
\usepackage{colortbl}
\usepackage{fancyhdr}
\usepackage{parskip}
\usepackage{ragged2e}
\usepackage{graphicx}
\usepackage{array}
\setlength{\parindent}{0pt}

\definecolor{tumblue}{RGB}{0,101,189}
\definecolor{tumgray}{RGB}{235,235,235}

\pagestyle{fancy}
\fancyhf{}
\renewcommand{\headrulewidth}{0pt}
\renewcommand{\footrulewidth}{0pt}
\fancyfoot[L]{\footnotesize \textbf{Technical University Munich} \\ Chair of Building Technology and Climate Responsive Design,\\School of Engineering and Design, Technical University of Munich}
\fancyfoot[R]{\footnotesize\thepage}

\fancyhf{}
\renewcommand{\headrulewidth}{0pt}
\renewcommand{\footrulewidth}{0pt}
\fancyhead[L]{\small\bfseries\color{tumblue}Technische Universit\"at M\"unchen}
\fancyhead[R]{\includegraphics[height=0.9cm]{LOGO_PATH}}
\fancyfoot[L]{\footnotesize \textbf{Technical University Munich} \\ Chair of Building Technology and Climate Responsive Design,\\School of Engineering and Design, Technical University of Munich}
\fancyfoot[R]{\footnotesize\thepage}

\fancypagestyle{firstpage}{%
  \fancyhf{}
  \renewcommand{\headrulewidth}{0pt}
  \renewcommand{\footrulewidth}{0pt}
  \fancyfoot[L]{\footnotesize \textbf{Technical University Munich} \\ Chair of Building Technology and Climate Responsive Design,\\School of Engineering and Design, Technical University of Munich}
  \fancyfoot[R]{\footnotesize\thepage}
}

\setlength{\tabcolsep}{2pt}
\setlength{\arrayrulewidth}{0.1pt}

\begin{document}
\thispagestyle{firstpage}


% --- Header -----------------------------------------------------------
\begin{minipage}[t]{0.7\textwidth}
{\LARGE\bfseries\color{black} COURSE_TITLE \\ SHEET_TITLE}
\end{minipage}%
\hfill
\begin{minipage}[t]{0.25\textwidth}
\raggedleft\includegraphics[width=1.5cm]{LOGO_PATH}
\end{minipage}

\bigskip
\noindent\begin{tabular}{@{}p{5cm}l@{}}
\textbf{Semester:} & SEMESTER \\
\textbf{Student:} & STUDENT_NAME \\
\textbf{Matriculation number:} & MATRICULATION_NUMBER \\
\textbf{Submission date:} & SUBMISSION_DATE \\
\end{tabular}

\bigskip
\hrule
\bigskip

% --- Feedback tables, totals, weighting, grading scale -----------------
FEEDBACK_TABLE_REPLACEMENT

\subsubsection*{Additional Comments}

Some more comments here.

\subsection*{Weighting}
In Assignment 2, most points focus on modeling, data use, and scenario analysis, as these reflect the main learning goals: developing a baseline, justifying modeling choices, and evaluating optimization scenarios. Formal requirements and presentation account for roughly 15\%. Clear formatting, figures, tables, and citations (APA) are important for professional practice and thesis preparation at TUM, and as climate engineers you will need to communicate data and models clearly to colleagues, clients, or regulatory bodies.

\subsection*{Grading Scale}
\begin{longtable}{|c|c|}
\hline
\textbf{Points Range} & \textbf{Grade} \\ \hline
> 76 & 1.0 - 1.3 \\ \hline
75 - 66 & 1.3 - 2.1 \\ \hline
65 - 56 & 2.1 - 2.7 \\ \hline
55 - 45 & 2.7 - 3.4 \\ \hline
44 - 37 & 3.4 - 4.0 \\ \hline
< 36 & 5.0 (fail) \\ \hline
\end{longtable}

\end{document}
